\chapter{总结与展望}
\section{主要工作与结论}
本文建立了CPU可运行、配置和原始结果可追溯的随机特征科研工程，验证Poisson五点差分、双子域解析导数及非零界面跳跃，并实现固定均匀、固定随机、PDE残差和联合指标采样。数学分析包括离散矩阵、局部尺度、加权最小二乘、标准右端扰动界及岭滤波解释。花形界面FD近似复现补充了与原文实现细节不同的诊断。

第一，合理静态基线不可缺少。当前平滑模型的联合策略没有稳定胜过均匀或随机点，固定更多界面点也可取得消融中的部分收益，不能只依据一个较差固定比例宣称方法改进。

第二，采样不能脱离表达力和目标解释。独立窄峰在当前固定特征及预算下未达到高精度，具体误差来源尚未唯一识别；联合策略在部分对照改善而未取得普遍优势。非均匀经验目标、配额和位置选择的效应需要区分，单纯增加点数的误差也不必单调。

第三，条件数、系数敏感性和物理解误差不是同一指标。尺度平衡可改善解却恶化物理通量；增宽可降误差同时增加病态；SVD截断会改变有效近似空间。固定矩阵的系数放大界有数值支持，但小函数变幅不足以证明一般稳定性。

第四，复现必须保留不一致。本文确认了原文花形数据和FD设置，却未取得足以逐值严格重现的全部细节，故明确标为近似复现。与原表更低或更高的误差都不自动构成新贡献。

\section{与最初研究目标的对应}
最低目标对应Poisson与简单界面的正确实现，已有模型一致性、单元测试和差分收敛数据支撑。标准目标对应采样对照、参数敏感性、误差与成本分析，已形成可审阅证据。拓展目标原本要求联合区域残差与界面误差的方法优于合理基线；当前证据没有达到这一条件。因此，研究方向与初衷一致，但不能把预期优势写成已实现的创新成果。

随机特征、局部网络、残差采样、SVD及界面约束均属于已有技术。本文可明确归于自身实施的工作是制造解与独立代码、共享特征的预算比较、同配额诊断、负面结果保留，以及有限维误差分解与扰动实验的组织。所给代数命题是标准工具的自含推导，不作为原创定理。尚未证明联合策略普遍优越，也未证明所有联合策略均无效；结论仅否定当前实验对强优势声明的支持。

\section{局限及后续方向}
研究仍局限于少量二维制造解、固定单隐层特征和五至十种子；正式版公式、RNN-PG全文和完整新颖性核查有缺口。实验没有建立随机自适应的全PDE误差理论，也没有完成学校官方模板确认。窄峰结果提示未来可研究更多局部子域、特征增长或算子感知采样，但这些已有相关工作\cite{ag,christoffel,schwarzpreprint}，应先明确差异，且把额外表达力和成本公平地给予对照。

后续优先处理独立监测池、尺度与重要性权重的关系、矩阵扰动及高精度抵消核查，并在新的模型上验证从已有结果形成的假设。论文的合理定位仍为既有方法的独立实现、可重复数值比较和稳定性分析；是否形成算法改进，应由更强证据决定。
